% Generated draft: requires mathematical and source review.
\documentclass[11pt]{article}
\usepackage[margin=25mm]{geometry}
\usepackage{fontspec}
\setmainfont{Latin Modern Roman}
\usepackage{amsmath,amssymb,mathtools,unicode-math}
\setmathfont{Latin Modern Math}
\usepackage{enumitem,needspace,xurl,hyperref}
\setlength{\parindent}{0pt}
\setlength{\parskip}{5pt}
\setlength{\emergencystretch}{3em}
\newcommand{\N}{\mathbb{N}}
\newcommand{\Z}{\mathbb{Z}}
\newcommand{\Q}{\mathbb{Q}}
\newcommand{\R}{\mathbb{R}}
\newcommand{\C}{\mathbb{C}}
\newcommand{\F}{\mathbb{F}}
\newcommand{\A}{\mathbb{A}}
\newcommand{\PP}{\mathbb{P}}
\newcommand{\E}{\mathbb{E}}
\newcommand{\eps}{\varepsilon}
\newcommand{\dd}{\,\mathrm{d}}
\begin{document}
\section*{UnsolvedMath ID: 701760}
\textbf{Source problem code:} catalogue entry
\subsection{Definitions and mathematical statement}
Let \(p\) be a prime. A pro-\(p\) group is a topological group that is an inverse limit of finite \(p\)-groups. It is finitely generated if some finite subset generates a dense subgroup. A pro-\(p\) group is \(p\)-adic analytic if it admits the structure of a finite-dimensional \(p\)-adic Lie group compatible with its topology.

All subgroups generated below are understood to be topologically closed. For closed subgroups \(K,L\leq G\), write
\[
K^{p^r}=\overline{\langle x^{p^r}:x\in K\rangle},\qquad
[K,L]=\overline{\langle x^{-1}y^{-1}xy:x\in K,\ y\in L\rangle}.
\]
Use the following conventions for the five named filtrations:
\begin{enumerate}
\item[(a)] The \(p\)-power series is \(S_i=G^{p^i}\) for \(i\geq0\).
\item[(b)] The iterated \(p\)-power series is \(S_0=G\) and \(S_{i+1}=S_i^p\) for \(i\geq0\).
\item[(c)] The lower \(p\)-series is \(P_1(G)=G\) and
\[
P_{n+1}(G)=\overline{P_n(G)^p[P_n(G),G]},\qquad S_i=P_{i+1}(G).
\]
\item[(d)] The Frattini series is \(S_0=G\) and \(S_{i+1}=\Phi(S_i)\), where
\[
\Phi(K)=\overline{K^p[K,K]}
\]
is the Frattini subgroup of a pro-\(p\) group \(K\).
\item[(e)] The dimension subgroup series is taken to mean the Zassenhaus filtration. Define \(\gamma_1(G)=G\) and \(\gamma_{j+1}(G)=[\gamma_j(G),G]\), and put
\[
D_n(G)=\overline{\left\langle \gamma_j(G)^{p^k}:j\geq1,\ k\geq0,\ jp^k\geq n\right\rangle},\qquad S_i=D_{i+1}(G).
\]
\end{enumerate}

For an infinite finitely generated pro-\(p\) group \(G\) and a filtration \(\mathcal S=(S_i)_{i\geq0}\) as above, define the Hausdorff dimension of a closed subgroup \(H\leq G\), relative to \(\mathcal S\), by
\[
\operatorname{hdim}^{\mathcal S}_G(H)
=\liminf_{i\to\infty}
\frac{\log_p[HS_i:S_i]}{\log_p[G:S_i]},
\]
where the ratio is taken for sufficiently large \(i\) with \([G:S_i]>1\). Its Hausdorff spectrum is
\[
\operatorname{hspec}^{\mathcal S}(G)
=\{\operatorname{hdim}^{\mathcal S}_G(H):H\leq G\text{ is closed}\}.
\]
Finite Hausdorff spectrum means that this set has finitely many elements.

\textbf{Question.} For each of the five filtrations \(\mathcal S\) separately, is the following implication true for every prime \(p\) and every infinite finitely generated pro-\(p\) group \(G\)?
\[
\bigl|\operatorname{hspec}^{\mathcal S}(G)\bigr|<\infty
\quad\Longrightarrow\quad
G\text{ is }p\text{-adic analytic}.
\]
The supplied question also allows finite \(G\); such groups are already \(p\)-adic analytic of dimension zero, so restricting the displayed formulation to infinite \(G\) avoids a finite-group convention for the dimension formula without removing a case requiring an answer. The five implications are separate questions, each assuming finiteness for its indicated filtration.

\textit{Unresolved source qualification.} The supplied record names the filtrations and Hausdorff spectrum without defining them. This draft uses the conventions above, including the spectrum over all closed subgroups and the Zassenhaus interpretation of the dimension subgroup series. Agreement with the cited original problem requires human source review.
\subsection{Short English statement}
For each of five standard filtrations, ask whether a finitely generated pro-\(p\) group must be \(p\)-adic analytic whenever its closed subgroups have only finitely many Hausdorff dimensions relative to that filtration. Each filtration gives a separate question. The supplied registry labels the problem open; that status and the filtration conventions have not been independently verified.
\subsection{Sources}
\begin{itemize}
\item Local UnsolvedMath source registry, supplied complete record: permanent ID \(701760\), exact source problem code \texttt{catalogue entry}, title \emph{Kourovka Notebook Problem 20.102}. The supplied record labels the problem \texttt{open} and gives an update date of \texttt{2026-10-16}; these are supplied registry metadata, not an independently verified literature assessment.
\item \emph{The Kourovka Notebook (Unsolved Problems in Group Theory)}, original Problem 20.102, 20th issue (2022); the supplied background also cites the 21st issue edition (2026), \texttt{arXiv:1401.0300v48}. The supplied collection description names E. I. Khukhro and V. D. Mazurov as editors. This citation has not been independently consulted.
\item Attribution: the statement begins with the credit ``Y. Barnea, A. Shalev.'' The background and proposer field separately identify A. Thillaisundaram as proposer, with proposed year 2022. Both attributions are retained without inferring their precise relationship.
\end{itemize}
\end{document}
